\documentclass[12pt]{article}
\usepackage[margin=1in]{geometry}
\usepackage{hyperref}
\linespread{1.3}

\title{Referee Report\\[6pt]
\large ``Frequent Losers as Cover Bidders in Public Procurement''\\
\large Genicolo-Martins \& Furquim de Azevedo}
\author{Simulated Referee Report (V4 Self-Assessment)}
\date{March 2026}

\begin{document}
\maketitle

\section*{Summary}

This paper proposes ``frequent losers'' (FL)---firms that never win yet
participate in abnormally many procurement tenders---as a screening
marker for cover bidding. Using 4.5 million tender-items from S\~ao
Paulo's BEC platform (2009--2019), the authors identify 2,735 FL firms
via an IQR-based threshold and document that FL-present tenders have
6--9\% higher prices (OLS). External validation against CADE convictions,
a leave-one-out IV, Bajari--Ye bid-coordination tests, and a rich set of
robustness checks round out the analysis.

The V4 revision represents a substantial improvement over V3: the H1/H2
framing clarifies what the paper claims, the LOO IV provides a credible
(if imperfect) causal strategy, the CADE permutation test strengthens
external validation, and the FL definition robustness exercises are
welcome. Nevertheless, several issues warrant attention before
publication. I organize my comments into major, medium, and minor
categories.

% =========================================================================
\section*{Major Issues}

\subsection*{1. IV Exclusion Restriction Needs Stronger Defense}

The LOO instrument counts FL firms active at \emph{other} PBUs in the
same product group and year. The exclusion restriction assumes that FL
supply at distant PBUs affects focal-PBU prices \emph{only} through FL
presence. This is plausible but not airtight. At least two threats
deserve discussion:

\begin{enumerate}[]
  \item \textbf{Correlated demand shocks within product-group--year cells.}
    If a statewide pharmaceutical shortage simultaneously raises prices
    across all PBUs and attracts FL firms (who specialize in specific
    product groups), the instrument would be correlated with unobserved
    price shocks. The item $\times$ year FE absorb aggregate item-level
    shocks, but the instrument varies at the item-group $\times$ PBU
    $\times$ year level, so within-item-type geographic shocks could
    violate exclusion.

  \item \textbf{Common supply-side shocks.}
    If a new FL firm enters a product market across multiple PBUs
    simultaneously (e.g., a newly created shell company bidding
    statewide), the instrument captures this firm's entry at other PBUs,
    but the \emph{same firm} may also enter the focal PBU for the same
    unobserved reason that affects prices.
\end{enumerate}

\noindent\textbf{Suggestion:} Add a paragraph explicitly discussing
these threats and why item + year + PBU FE mitigate them. Consider
testing robustness to narrower geographic definitions of the LOO set
(e.g., only PBUs in \emph{different} municipalities). If the IV estimate
is stable across geographic radii, this strengthens the exclusion
restriction. Also discuss the Goldsmith-Pinkham et al.\ (2020)
Bartik-style decomposition more concretely---which ``shares'' drive the
instrument, and are they plausibly exogenous?

\subsection*{2. The Network Split Tells a Different Story Than Originally Hypothesized}

The paper classifies FL firms by winner HHI and repeat-partner counts,
labeling them ``high-concentration'' and ``low-concentration.'' The
original design (evident from the variable names \texttt{has\_high\_susp\_fl}
and \texttt{has\_low\_susp\_fl}) was to separate genuine cover bidders
(``high-suspicion'') from incompetent firms (``low-suspicion''). The
expectation was that high-suspicion FL firms would drive the price effect.

The results are the \emph{opposite}: low-concentration (originally
``low-suspicion'') FL firms drive the entire effect (0.126, $p < 0.01$),
while high-concentration FL firms show no significant effect ($-0.018$).
The paper reframes this as ``cover bidding is most profitable in
competitive markets,'' which is economically sensible but represents a
post-hoc rationalization.

\noindent\textbf{Concerns:}
\begin{enumerate}[]
  \item The classification was designed to identify cover bidders, not to
    test market-structure heterogeneity. If the goal is market-structure
    heterogeneity, why not simply interact FL presence with market HHI
    directly? The current classification conflates firm-level network
    characteristics with market-level concentration.

  \item The negative (though insignificant) coefficient on
    high-concentration FL firms undermines the claim that FL firms are
    cover bidders. If 1,356 of 2,735 FL firms show no price effect, are
    they genuinely incompetent firms that should not be flagged by the
    screen? This has direct implications for the welfare calculation.

  \item The renaming from ``high-suspicion'' to ``high-concentration''
    without changing the underlying analysis is cosmetic. The reader
    should be told that the original hypothesis was not confirmed and
    that the reinterpretation, while plausible, was arrived at ex post.
\end{enumerate}

\noindent\textbf{Suggestion:} (i) Present a direct interaction of FL
presence with market-level HHI (continuous), independent of the network
classification. (ii) Acknowledge transparently that the network
classification did not perform as originally designed. (iii) Recalculate
welfare losses using only the low-concentration FL subset, as these are
the firms driving the price effect.

\subsection*{3. Bajari--Ye Fake-Groups Placebo Failure}

The fake-groups placebo is designed to validate the Bajari--Ye
methodology: randomly splitting non-FL firms into two groups should
yield null results. Instead, the placebo pairwise product is 5.80
(\emph{larger} than the FL product of 5.16) and statistically
significant.

The paper argues that this reflects ``common tender-level shocks'' and
that the FL vs.\ non-FL comparison is the relevant test. However:

\begin{enumerate}[]
  \item The FL pairwise product (5.16) is \emph{lower} than the
    fake-group product (5.80). Under the cover bidding hypothesis, FL
    pairs should exhibit \emph{stronger} correlation than random
    non-FL pairs, not weaker. The paper's argument that the FL--non-FL
    difference (2.95) is what matters does not address why FL pairs are
    \emph{less} correlated than random non-FL pairs.

  \item The magnitude ordering (non-FL: 2.21 $<$ FL: 5.16 $<$
    fake-groups: 5.80) is puzzling. If common tender-level shocks drive
    positive pairwise products, one would expect FL and non-FL products
    to be similar (both affected by shocks), with the FL product
    \emph{additionally} elevated by coordination. Instead, the non-FL
    product is much lower, suggesting that the FL/non-FL partition itself
    captures some feature of the tender environment (e.g., FL-present
    tenders are in more volatile markets), not necessarily bid
    coordination.

  \item The ``fake groups'' placebo was designed by the authors as a
    validity check. Its failure should be discussed more honestly as a
    limitation of the Bajari--Ye implementation, rather than reframed as
    confirmation.
\end{enumerate}

\noindent\textbf{Suggestion:} (i) Compute the fake-groups placebo
separately for FL-present and FL-absent tenders. If common shocks are
the explanation, both subsamples should produce similar placebo products.
(ii) Report the within-tender variance of bid residuals for FL vs.\
non-FL tenders to diagnose whether FL-present tenders are systematically
noisier. (iii) Consider alternative residualization strategies (e.g.,
tender FE in the first-stage bid equation) that absorb common shocks
before computing pairwise products.


\subsection*{4. Null DiD Result and Causal Inference}

The Callaway \& Sant'Anna estimator yields ATT = 0.014 (SE = 0.039),
statistically insignificant. This is the one design that exploits
within-market temporal variation---arguably the cleanest identification
strategy---and it produces a null result. The paper acknowledges this
and frames DiD as ``supplementary evidence,'' arguing that limited power
from gradual FL entry prevents detection.

While the power argument has merit, a skeptical reader could note:

\begin{enumerate}[]
  \item The null DiD is also consistent with no causal effect: FL firms
    select into expensive markets, and the cross-sectional association
    reflects this selection rather than a causal mechanism.

  \item The sample (19,777 markets, 1,511 treated) is not trivially
    small. The standard error (0.039) would detect an effect of 0.077 or
    larger at the 5\% level. Since the OLS estimate is 0.064, the DiD
    does have \emph{some} power to detect effects of the claimed
    magnitude.

  \item The Rambachan--Roth bounds are described as ``wide'' but no
    specific numbers are reported in the main text. What are the honest
    confidence sets under $\bar{M} = 0$, 0.01, and 0.02?
\end{enumerate}

\noindent\textbf{Suggestion:} (i) Report the minimum detectable effect
(MDE) for the DiD design explicitly, so readers can assess power. (ii)
Report the Rambachan--Roth honest confidence sets for specific values of
$\bar{M}$ in the robustness table. (iii) Discuss more candidly that the
null DiD weakens the causal claim in H2, rather than treating it as
uninformative.


% =========================================================================
\section*{Medium Issues}

\subsection*{5. CADE Validation: Modest Unconditional Effect}

The unconditional co-participation odds ratio is 1.21 ($p = 0.017$).
This is statistically significant but economically modest: FL firms are
only 21\% more likely than non-FL firms to co-participate with CADE
cartelists. The 3.5$\times$ figure comes from the permutation test,
which controls for participation intensity by matching on participation
quartiles.

The permutation-adjusted comparison is more compelling, but the paper
leads with the 3.5$\times$ figure (abstract, introduction) without
making clear that the unconditional OR is only 1.21. This creates an
impression of stronger external validation than the raw data support.

\noindent\textbf{Suggestion:} Report both figures prominently---the
unconditional OR (1.21) and the participation-adjusted ratio
(3.5$\times$)---in the abstract and introduction, so the reader can
judge the strength of validation.

\subsection*{6. Cross-Fitting Attenuation}

The cross-fitting exercise yields an average coefficient of 0.036,
attenuated 44\% relative to the full-sample OLS (0.064). One fold is
significant (0.053, $p < 0.05$), the other is not (0.019, $p > 0.10$).
The paper attributes the attenuation to ``lower statistical power from
fewer FL firms identified in even years.''

However, cross-fitting is designed precisely to detect overfitting---the
attenuation could reflect the full-sample estimate being inflated by
in-sample FL definition. If the ``true'' effect is closer to 0.036 than
0.064, the welfare estimates should be adjusted accordingly.

\noindent\textbf{Suggestion:} (i) Report the cross-fit coefficient in
the welfare analysis as a third bound (between OLS and IV). (ii)
Discuss why one fold is significant and the other is not beyond ``power
differences''---are there structural differences between odd and even
years (e.g., election cycles, policy changes)?

\subsection*{7. Homogeneous Sub-Samples Infeasible}

Table B.5 reports ``Analysis infeasible'' for the homogeneous sub-sample
analysis (fuel, A4 paper, generic medicines). The explanation is that
BEC item codes do not map to product-level descriptions. This is a
notable gap: the homogeneous-goods test was specifically requested to
rule out quality-selection stories, and the inability to conduct it
leaves this alternative explanation partially unaddressed.

\noindent\textbf{Suggestion:} (i) Even without perfect product labels,
some BEC item codes likely correspond to clearly homogeneous goods.
Consider a manual inspection of the top-50 most common item codes to
identify commodities. (ii) As a partial substitute, report the FL
coefficient separately for item groups with low vs.\ high price
dispersion (within-item-group CV of prices), on the logic that
low-dispersion items are more homogeneous.

\subsection*{8. Welfare Range Too Wide}

The welfare estimates span R\$734 million (OLS) to R\$2.4 billion (IV),
a 3.3$\times$ range. This is driven entirely by the OLS--IV coefficient
gap (0.064 vs.\ 0.194). While wide bounds are honest, they limit policy
relevance. Moreover, if the exclusion restriction is imperfect (Issue
\#1), the IV upper bound may be inflated.

\noindent\textbf{Suggestion:} (i) Add a ``preferred'' welfare estimate
using the cross-fit coefficient (0.036) or the network-split coefficient
for low-concentration FL firms (0.126). (ii) Compute the welfare loss
attributable to FL firms in low-concentration markets only, since these
drive the price effect. (iii) Conduct a bounding exercise: if the IV is
biased upward by $X$\%, what is the implied welfare loss?

\subsection*{9. Convite vs.\ Preg\~ao Tension with Oversight Hypothesis}

The oversight heterogeneity analysis shows the FL coefficient is larger
for smaller PBUs (less oversight $\rightarrow$ more cover bidding),
supporting Prediction~5. However, the main results show that preg\~ao
(electronic, more transparent, typically larger PBUs) has a \emph{larger}
FL coefficient (9.3\%) than convite (sealed-bid, less transparent, 3.8\%).

The paper explains this as preg\~ao's real-time bid observation
facilitating Regime~1 cover bidding. While plausible, this tension
between the oversight hypothesis (more oversight $\rightarrow$ smaller
effect) and the modality result (more transparent modality
$\rightarrow$ larger effect) should be discussed more explicitly. The
two results pull in opposite directions regarding the role of
transparency.

\noindent\textbf{Suggestion:} Add a paragraph in Section~8 or
Section~9 explicitly acknowledging this tension and discussing whether
``oversight'' and ``transparency'' operate through different channels.

\subsection*{10. Placebo IV Interpretation}

The placebo IV uses sub-threshold always-losers as the instrument source
and yields a coefficient of 5.24 with a very weak first stage (coef =
0.00004). The paper correctly identifies this as a weak-instrument
artifact. However:

\begin{enumerate}[]
  \item A true placebo should produce a \emph{null} result, not a
    wild artifact. The fact that the sub-threshold instrument has
    \emph{any} first-stage correlation (even tiny) with FL presence
    raises questions about whether the main instrument's first stage
    captures FL-specific variation or broader always-loser market
    dynamics.

  \item The correct placebo IV test would instrument FL presence with
    sub-threshold always-loser supply and check whether the 2SLS
    estimate is zero. Instead, the enormous coefficient signals
    near-zero relevance, which is uninformative about the exclusion
    restriction.
\end{enumerate}

\noindent\textbf{Suggestion:} (i) Consider a more informative placebo:
instrument FL presence with the supply of firms that \emph{do} win
(win rate $> 0$) at other PBUs. These firms should not affect prices
through FL mechanisms. (ii) Report Anderson--Rubin confidence sets for
the placebo IV, which are robust to weak instruments.


% =========================================================================
\section*{Minor Issues}

\subsection*{11. Abstract Length}

The abstract is approximately 140 words, which is appropriate. However,
the welfare figures (R\$734M to R\$2.4B) take up space without adding
much in an abstract. Consider replacing with a qualitative statement
(``substantial welfare losses'') to save space for the IV result and
CADE validation, which are more novel.

\subsection*{12. Prediction~3 (Regime Identification) Lacks Clear Test}

Prediction~3 states that Regime~1 predicts higher FL dispersion and
Regime~2 predicts lower FL dispersion. The paper concludes ``the
empirical pattern most closely resembles Regime~1.'' However, no formal
test is reported. What is the dispersion ratio? Is the difference
statistically significant? The regime density comparison (Figure~7) is
visual only.

\noindent\textbf{Suggestion:} Add a formal variance-ratio test or
Kolmogorov--Smirnov test comparing FL and non-FL bid spread
distributions, with a p-value.

\subsection*{13. Table Formatting Inconsistencies}

Several appendix tables use different formatting conventions (some have
\texttt{adjustbox}, some don't; some include \texttt{threeparttable}
notes, some use captions only). Standardize across all tables for
consistency.

\subsection*{14. Missing Figures in Appendix}

Several figures referenced in the plan are absent from the appendix
(e.g., bootstrap histogram, HonestDiD overlay, cover bid spread
density). While the paper notes some are ``omitted,'' the reader would
benefit from knowing which figures were attempted and why they were
dropped. Add brief comments.

\subsection*{15. ``Conceptual Framework'' vs.\ Formal Model}

The conceptual framework introduces formal notation (FOC, definitions,
predictions) but stops short of deriving equilibrium. This is fine given
the paper's empirical focus, but the comparative statics (``$m^*$
increasing in $n$ and decreasing in $\theta$'') are stated without
proof. Since the model is simple enough, either show the derivation in
a short appendix or cite a paper that does.

\subsection*{16. Temporal FL Definition}

The temporal FL (3-year rolling windows) produces a coefficient of 0.070,
\emph{larger} than the full-sample 0.064. This is unexpected---if
anything, temporal FL should introduce noise (shorter windows, fewer
observations per firm) and attenuate the coefficient. The result
deserves brief discussion: does the temporal definition capture different
firms at different times? How much overlap is there between the
full-sample and temporal FL sets?

\subsection*{17. Notation}

Equation~(1) uses $\alpha_g$, $\lambda_t$, $\gamma_k$ for item, year,
and PBU FE, but the subscripts $g$, $t$, $k$ are not consistently
defined in the surrounding text. Clarify that $g$ indexes item groups
(not individual items), $t$ indexes years, and $k$ indexes PBUs.

\subsection*{18. BEC Representativeness}

The paper analyzes only S\~ao Paulo's BEC platform. How representative
is BEC of Brazilian procurement more broadly? If BEC has above-average
transparency (it is electronic), the FL screen's performance may differ
in less transparent states. A brief discussion of external validity
would strengthen the conclusion.


% =========================================================================
\section*{Overall Assessment}

\paragraph{Strengths.}
\begin{itemize}[]
  \item The H1/H2 framing is clear and honest about what the paper
    establishes (H1) vs.\ what it suggests (H2).
  \item The FL screen is operationally practical---it requires only
    participation and outcome data, making it deployable in
    data-poor environments.
  \item The robustness portfolio is extensive: threshold sensitivity,
    low-win-rate variants, cross-fitting, temporal FL, matching,
    sensitivity analysis, and alternative clustering.
  \item The LOO IV strategy is well-motivated and the first stage is
    strong ($F = 396$).
  \item The CADE permutation test with stratified resampling is a
    thoughtful validation exercise.
  \item The paper is clearly written, well-organized, and the 50-page
    length is appropriate for the scope.
\end{itemize}

\paragraph{Weaknesses.}
\begin{itemize}[]
  \item The IV exclusion restriction is not airtight and deserves more
    thorough discussion of threats.
  \item The network split results contradict the original hypothesis
    and the reinterpretation, while plausible, is post hoc.
  \item The Bajari--Ye fake-groups placebo fails, undermining
    the bid-coordination evidence.
  \item The null DiD result weakens the causal claim, and the power
    argument should be quantified.
  \item The unconditional CADE OR of 1.21 is modest; the 3.5$\times$
    figure requires participation adjustment that should be made
    more transparent.
  \item The cross-fitting attenuation (0.036 vs.\ 0.064) is
    substantial and may indicate some overfitting in the FL definition.
  \item The welfare range (R\$734M--R\$2.4B) is too wide for actionable
    policy guidance.
\end{itemize}

\paragraph{Recommendation.}
The paper makes a genuine contribution by proposing a practical cartel
screening tool and providing multi-pronged evidence for its validity.
The V4 revision substantially strengthens the analysis relative to V3.
However, the issues above---particularly the IV exclusion restriction
defense, the honest treatment of the network split and Bajari--Ye
placebo results, and the null DiD---should be addressed before
publication. I recommend \textbf{major revision} with the understanding
that the revisions are primarily about \emph{presentation and
interpretation} rather than new analyses, and should be achievable in
a single round.

\end{document}
